\section{Exceptional sheaves at a symmetric quantum parameter}\label{sec:surfaces}

For a del Pezzo surface there is a way to recognize the individual objects of an asymptotic basis. An exceptional object on the surface has just one nonzero cohomology sheaf, and the integral class determines that sheaf. This rigidity is especially effective when two idempotents have the same Euler value. The Stokes matrix sees only their Euler pairing. On these surfaces one can recover the vanishing of the entire morphism complex.

\begin{theorem}\label{thm:surface}
Let $X$ be a smooth complex del Pezzo surface and $G$ a finite group of automorphisms. Let $\tau$ be a $G$-fixed semisimple point in the convergent domain of Gamma conjecture II with its $G$-compatible descendant calibration. For each admissible lifted phase there is a full exceptional collection of sheaves $(F_s)_{s\in S}$, indexed by the primitive idempotents at $\tau$, with the following properties.
\begin{enumerate}
\item The calibrated Gamma section of $F_s$ is the normalized asymptotic column $Y_s$ for a choice of signs on the idempotent orbits. Its labelled line determines $F_s$ up to isomorphism, and $gF_s\simeq F_{gs}$.
\item The collection is ordered by the projected Euler values. If $s\ne t$ and $u_s=u_t$, then $\mathrm{RHom}(F_s,F_t)=0$ in both directions.
\end{enumerate}
Such a fixed parameter exists for every toric del Pezzo surface and every finite group of its automorphisms.
\end{theorem}

\begin{proof}
Use the orbitwise normalization and the common continuation in $z$ of Section~\ref{sec:lattice}. Hu, Ke, Li and Su prove the ordinary Gamma conjecture II for del Pezzo surfaces \cite[Theorem~1.4 and \S2.3]{HKLS}. Applied to our labelled columns, it gives a full exceptional sequence $(A_s)$ in the required weak phase order, with $\mathfrak{s}_\tau(A_s)=c_sY_s$ for some scalars $c_s\ne0$. Only the lines of their sections enter here. To fix the scalars in our normalization, pair each section with itself using the flat bilinear pairing of \cite[\S3.4]{GGI}. The self-pairing on the left is $\chi(A_s,A_s)=1$ by Riemann--Roch \cite[Lemma~3.4.6]{GGI}. The self-pairing of $Y_s$ is the diagonal Stokes entry $1$ \cite[Proposition~2.6.4]{GGI}. It follows that $c_s^2=1$. Shifting $A_s$ by one if necessary gives the initially chosen sign for every $s$.

The classes of a full exceptional sequence generate $K_0(X)$, and its unitriangular Euler matrix makes them a basis. Riemann--Roch therefore gives injectivity of the Chern character on this lattice. Naturality of the framing and uniqueness of the asymptotic solution now give $[gA_s]=[A_{gs}]$, as in the proof of Proposition~\ref{thm:main}.

By \cite[Lemma~4.11]{Krah}, exceptional objects on $X$ with the same class differ by an even shift. Each $A_s$ has the form $F_s[n_s]$ for an exceptional sheaf $F_s$ \cite[Propositions~2.9--2.10]{KO}. Pullback by an automorphism is exact on coherent sheaves. Comparing the single nonzero cohomology sheaves in $gA_s\simeq A_{gs}[2k]$ gives
\[
               gF_s\simeq F_{gs},\qquad n_s\equiv n_{gs}\pmod 2 .
\]
Removing the shifts thus changes signs uniformly on each orbit. Shifts preserve the generated subcategory and semiorthogonality. For prescribed signs one may instead use $F_s[n_s\bmod 2]$. The integral points on each labelled line are generated by a primitive class since its Euler self-pairing is one. The sheaf representative singles out one sign: an exceptional bundle has positive rank, while an exceptional torsion sheaf is supported on a $(-1)$-curve and its first Chern class has positive anticanonical degree. Equality of classes between exceptional sheaves gives an isomorphism by the same rigidity.

For distinct $s,t$ with $u_s=u_t$ the off-diagonal Stokes entries vanish \cite[Proposition~2.6.4]{GGI}. Riemann--Roch identifies the corresponding Euler pairings with zero. In an exceptional pair of coherent sheaves on $X$, at most one forward Ext group is nonzero \cite[Corollary~2.11]{KO}. Its alternating dimension can vanish only if every such group vanishes. Backward vanishing already follows from exceptionality. The argument applies even when $s,t$ lie in different $G$-orbits.

For the final assertion let $W$ be the full fan automorphism group. Its map to $\mathrm{Aut}(X)/\mathrm{Aut}^0(X)$ is split surjective \cite[Theorem~4.5]{Dun}. The connected group acts trivially on cohomology, since each of its elements is homotopic to the identity. Lift the component image of $G$ to a subgroup of $W$ and apply Proposition~\ref{thm:main}. Its fixed parameter is therefore fixed by $G$. The stable-map construction gives the required naturality for the geometric action of $G$ itself.
\end{proof}

\subsection{Restriction to a subgroup}
Near $\tau$ the canonical coordinates $u_s$ remain independent even when their values at $\tau$ coincide. The group permutes the coordinates with their labels, so for a subgroup $J\subset G$ the fixed slice has one coordinate for each $J$-orbit. On a small ball the normalized sectorial solution extends as a flat matrix in the parameters \cite[Proposition~2.5.1]{GGI}. After continuation within this ball the labelled Gamma lines are the same wherever the chosen phase is admissible. By Theorem~\ref{thm:surface} the sheaves are then the same. At a generic point of the $J$-fixed slice the old equal-value blocks split into $J$-orbits. Objects whose order changes within an old block are mutually orthogonal and can be interchanged. Thus restriction has a literal local meaning for these subcategories. A continuation across a genuine simple Stokes wall can mutate the collection \cite[\S4.6]{GGI}.

For instance, the mirror of the degree-six toric del Pezzo at equal ray coefficients has values $6t,-3t,-2t$ for a common positive scale $t$. The orbit cardinalities at these values are respectively one, two and three. Indeed in hexagonal coordinates the six monomials are $x,y,xy$ and their inverses. The point $(1,1)$ gives $6t$, the two diagonal nontrivial cube roots give $-3t$, and the three nonidentity sign pairs give $-2t$. These are all the critical points: in the chosen coordinates the logarithmic derivative equations are
\[
 x-x^{-1}=y-y^{-1}= (xy)^{-1}-xy .
\]
The first equality gives $x=y$ or $xy=-1$. Substitution on the diagonal gives $(x^2-1)(x^2+x+1)=0$, and substitution off it gives the two remaining sign pairs. To see nondegeneracy, the logarithmic Hessian divided by $t$ has diagonal entries $x+x^{-1}+a,\ y+y^{-1}+a$ and off-diagonal entries $a$, with $a=xy+(xy)^{-1}$. Its determinants on these three sets are $12,3,-4$, respectively. Hexagon symmetries permute the three sign pairs transitively and inversion exchanges the two cube roots. The theorem gives three blocks with these labels and total orthogonality inside each block at any sufficiently small $t>0$. Previously constructed equivariant surface collections give a geometric description of this example \cite{BDM,ESS}.

On the $m$-th power of the degree-six surface the full fan-group orbits are indexed by the numbers $(a,b,c)$ of entries from these three surface orbits. Their sum is $m$. The orbit has size $m!2^b3^c/(a!b!c!)$ and its Euler value divided by $t$ is $-2m+8a-b$. Two different orbits can therefore have the same value only if their occupation numbers differ by a nonzero multiple of $(1,8,-9)$. The first coincidence occurs for $m=9$. A tuple from $(1,8,0)$ and one from $(0,0,9)$ both have value $-18t$. One orbit has $2304$ members and the other $19683$. Proposition~\ref{prop:product} gives vanishing in every degree between any two distinct sheaves in their union. The three orbit values $u_1,u_2,u_3$ on the surface are independent local coordinates of the fixed germ. A small generic change breaks $u_1+8u_2=9u_3$, splitting this block into two while keeping the same external sheaves at a nearby admissible parameter for the chosen phase.

\subsection{The degree-five surface and its three atoms}\label{sec:quintic}
The quintic del Pezzo surface $X$ is the blow-up of the plane at four points in general position. Its five conic pencils give a faithful action of $\mathrm{Aut}(X)=S_5$ on a set of geometric objects. Write $h_1,\ldots,h_5$ for their fibre classes. Four of these are $H-E_i$, and the fifth is $2H-\sum_i E_i$ in the usual blow-up basis. Each has square zero, and two different ones meet once. Their Gram matrix is invertible, so $H^2(X,\C)$ is the five-letter permutation representation. The primitive part $c_1^\perp$ is its standard four-dimensional summand.

On this surface an explicit equivariant atomic decomposition is known \cite[Lemma~3.15 and Proposition~5.3]{ESS}, with an earlier collection on $\overline M_{0,5}$ in \cite[Example~1.3]{CT}. Castravet and Tevelev explicitly associate the two simple eigenvalues to the structure sheaf and a rank-two bundle, and the eigenvalue $-3$ to five line bundles \cite[Introduction]{CT}. Their comparison motivates the question of recognizing the equivariant components of the calibrated collection at each phase. The standard geometric decomposition is generated by a rank-two exceptional bundle, the five line bundles $\mathcal O(-h_i)$, and the structure sheaf. Each singleton generates a copy of the derived category of a point. The five line bundles generate a sum of five such categories with the transitive $S_5$-action. Here the word \emph{atom} refers to these equivariant admissible components, taken up to the whole-block mutations of \cite{ESS}. We can recognize them using only the asymptotic columns and the phase order.

\begin{proposition}\label{prop:quintic}
At the small divisor parameter $(\log t)c_1(X)$ for sufficiently small $t>0$, the primitive idempotents have three Euler values. Two are singletons with values $5t(1\pm\sqrt5)/2$. The five at $-3t$ form the $S_5$-set of conic pencils. At every admissible lifted phase the three subcategories generated by the exceptional sheaves of Theorem~\ref{thm:surface} give the above atomic decomposition up to $S_5$-equivariant whole-block mutations. Stabilizers act with trivial multiplier on the simple generators of these components.
\end{proposition}

\begin{proof}
The divisor equation weights a curve of class $\beta$ by $t^{c_1\cdot\beta}$. In the product of classes of half-degrees $i,j$, the coefficient in half-degree $l$ has weight $t^{i+j-l}$. Multiplying $H^{2i}$ by $t^i$ therefore identifies the small algebra at $t$ with that at $1$, and takes the Euler element to $tc_1$. The isomorphism commutes with $S_5$. Small positive $t$ lies in the large-radius convergence domain \cite[Remark~2.1]{HKLS}.

At $t=1$ Bayer and Manin compute multiplication by $c_1$ \cite[Table~3.5.3 and Proposition~3.5.4]{BM}. On the span of $1,c_1,[\mathrm{pt}]$ its characteristic polynomial is
\[
                    (\lambda+3)(\lambda^2-5\lambda-25);
\]
on $c_1^\perp$ the operator is $-3$. It is diagonalizable, so its three spectral projections split the algebra into ideals. The only issue is the radical of the five-dimensional ideal $A$ at $-3$. As a representation, $A$ consists of its unit line and the faithful standard summand $V$. The radical is invariant and cannot contain the unit. If nonzero, it must be $V$, and its square must be zero by nilpotence and irreducibility. The Frobenius pairing restricts nondegenerately to $A$ since distinct factors are orthogonal. A square-zero $V$ would be totally isotropic, which is impossible for a four-dimensional subspace of a five-dimensional nondegenerate space. All seven factors are consequently fields. Semisimplicity also occurs in Sanda's Floer-theoretic calculation \cite[\S6.2]{Sanda}. Smith gives the full ring independently \cite[\S4.3]{Smith}.

In $A$ there is one orbit of primitive idempotents, since its invariant subspace has dimension one. Its faithful action on five letters is the usual one up to isomorphism, since all automorphisms of $S_5$ are inner. The stabilizer $S_4$ is self-normalizing, so there is a unique equivariant bijection of this orbit with the conic pencils. This identifies labels before making any choice of mutation.

By Theorem~\ref{thm:surface} the sheaves form a full three-block collection of type $\{1,1,5\}$. All are bundles by \cite[Corollary~3.4]{KN}. The table of Karpov and Nogin \cite[\S5.7, row (4)]{KN} gives precisely one orbit of such collections under block mutations and simultaneous line-bundle twists, with arbitrary reorderings inside blocks. Starting with the displayed geometric collection, perform the mutations first and leave the common twist by $L$ until the end. Mutations through invariant admissible components give equivariant equivalences on the components by the natural projection functors \cite[Remark~3.23 and Lemma~3.24]{ESS}. Passing to the usual sheaf representatives only shifts generators. A singleton generator $B$ just before the twist and its image $B\otimes L$ in the marked sheaf collection are both fixed by $S_5$. Indeed the only exceptional generators of the category of a point are shifts of one another, and pullback preserves sheaves. Taking determinants gives
\[
                 \operatorname{rk}(B)(g^*L-L)=0
                        \qquad\text{in }\operatorname{Pic}(X).
\]
The Picard group is torsion-free. Its invariant line is generated integrally by $K_X$, which is primitive in the blow-up basis. Thus $L$ is a power of $\omega_X$. The helix rotation of a three-block collection moves the first block to the end with an anticanonical twist \cite[\S5.6]{KN}. Three rotations implement the simultaneous twist, so a further block mutation word absorbs $L$.

Finally the standard components have zero stabilizer multiplier by \cite[Proposition~5.3]{ESS}. This means that the projective action on a chosen generator can be made linear for its stabilizer, $S_5$ for a singleton and $S_4$ for a member of the five-block. Equivariant equivalences carry these structures through the mutations.
\end{proof}

A phase and a calibration specify the individual sheaves in Theorem~\ref{thm:surface}. Proposition~\ref{prop:quintic} identifies the equivariant mutation class of the resulting decomposition. Identifying its five sheaves with the five displayed line bundles at a prescribed phase would additionally require the corresponding five integral classes in the blow-up basis.

\subsection{A boundary of the fixed-locus theorem}\label{sec:weyl}
The existence assertion for toric varieties has a sharp boundary among del Pezzo surfaces with invariant Picard rank one. The distinction is visible in the algebra before choosing a phase: in low degrees an entire ideal of nilpotents survives all symmetry-preserving deformations. We use the usual logarithmic divisor coordinates on the branch from large radius.

\begin{proposition}\label{prop:cutoff}
Let $X$ be a smooth complex del Pezzo surface with a finite automorphism group $G$ acting with $\operatorname{rk}\operatorname{Pic}(X)^G=1$. Its logarithmic fixed locus meets the semisimple locus on the large-radius branch if and only if $d=K_X^2\ge5$. In degrees $d\le4$ write $r=9-d$ and $V=K_X^\perp\subset H^2(X,\C)$. Throughout the convergent fixed germ the quantum algebra has an ideal
\[
                 I=V\oplus\C\gamma,\qquad I^2=\C\gamma,\qquad I^3=0 .
\]
Generically on the component containing the small positive anticanonical ray, the algebra is a product of $\C$ and a local Frobenius algebra of length $r+2$ with maximal ideal $I$.
\end{proposition}

The cutoff at the monotone small parameter is known \cite[\S1.2]{Kaw}. Bayer and Manin describe the low-degree algebra at that parameter and compute the scalar action of the invariant classes also with point insertions \cite[Corollary~3.3.2 and Proposition~3.5.5]{BM}. We give the deduction for the entire fixed germ because it explains which deformations preserve the obstruction. Elagin, Schneider and Shinder discuss the corresponding change in the number of surface atoms \cite[Remark~1.6]{ESS}. For a cubic, Cavenaghi, Katzarkov and Kontsevich describe the two generic Euler values on this fixed locus \cite[Example~2.14]{CKK}.

\begin{proof}
In degree five Proposition~\ref{prop:quintic} works for the full automorphism group. In degrees at least six every smooth complex del Pezzo surface is toric, including both surfaces of degree eight. Theorem~\ref{thm:surface} supplies the positive assertion.

For $d\le4$ the condition on divisors gives a fixed parameter in the span $T$ of $1,c_1,[\mathrm{pt}]$. The entire Weyl group $W$ of the root lattice $K_X^\perp$ acts on the quantum algebra there. Indeed the point correlators are invariant under Weyl relabelling of the curve classes \cite[\S5.1]{GP}. The divisor equation and the unit axiom then give the covariance of the full product \cite[(3.1)--(3.3) and \S3.3]{BM}. The divisor weights are invariant on the chosen logarithmic slice, and the point parameter is arbitrary in the convergence germ. Thus these identities of series hold throughout that germ. The argument also applies after continuation preserving this Weyl action.

As a $W$-representation the algebra is $T\oplus V$, where $T$ is trivial and $V$ is the irreducible reflection representation. An invariant element $a\in T$ acts on $V$ by a scalar $\chi(a)$. This gives a unital algebra character $\chi$ of $T$. For $v,w,x\in V$, the trilinear form $(v*w,x)$ is symmetric and invariant. None of the reflection groups $D_5,E_6,E_7,E_8$ has a homogeneous cubic invariant on $V$. Equivalently the symmetric cubic sum over each Weyl orbit vanishes \cite[Lemma~3.5.2]{BM}. The same invariant-theoretic obstruction occurs in \cite{ADE}. Consequently $v*w$ has zero component in $V$. Write $\gamma\in T$ for the vector dual to $\chi$ under the Poincar\'e pairing restricted to $T$. Frobenius invariance gives
\[
                (v*w,a)=(v,w*a)=(v,w)\chi(a),
              \qquad v*w=(v,w)\gamma .
\]
In particular $(\gamma,1)=1$. Choose $x,y\in V$ perpendicular with $(x,x)=1$ and $y\ne0$. Associativity gives $\gamma*y=x*(x*y)=0$, whence $\chi(\gamma)=0$. Also $a*\gamma=(a*x)*x=\chi(a)\gamma$, so $\gamma^2=0$. This proves all the assertions about $I$ and excludes semisimplicity.

The quotient by $I$ has dimension two. At the monotone small parameter Bayer and Manin's explicit product has one reduced factor and one local factor, whose maximal ideal has dimension $r+1$ \cite[Proposition~3.5.5]{BM}. The dimension there identifies the radical with $I$. The grading change in the proof of Proposition~\ref{prop:quintic} applies in every degree, so the quotient is split semisimple also on the small positive ray. Nondegeneracy of the trace pairing is open, and its determinant is analytic. Hence the quotient is generically split semisimple on the indicated component. Its two idempotents lift uniquely across the nilpotent ideal. They are $W$-invariant since the quotient is a quotient of $T$. Exactly one of them has $\chi$-value one. This idempotent acts as the identity on $I$, and the other kills $I$, giving the stated lengths.
\end{proof}

The logarithmic lift in Proposition~\ref{prop:cutoff} specifies which parameters are fixed. Fixed points after quotienting divisor coordinates by integral periods may lie on other components. In some examples an obstruction can be read from the geometric group alone. The deck involution of a degree-two double plane has trace $-4$ on total cohomology \cite[Example~2.5]{CT}. At any fixed semisimple quantum parameter with this action its trace would count the primitive idempotents it fixes. The negative trace excludes every such parameter.
